\documentclass{ubcthesis}
\usepackage[utf8]{inputenc}
\usepackage[T1]{fontenc}
\usepackage{geometry}
\usepackage{graphicx}
\usepackage{float}
\usepackage{hyperref}
\usepackage[english,indonesian]{babel}

\title{Optimasi Pemrosesan Template LaTeX Menggunakan Pendekatan Berbasis AI}
\author{Academic Support Agent}
\date{\today}

\begin{document}
\maketitle

\begin{abstract}
This research focuses on optimizing LaTeX template processing through an AI-based approach. The study demonstrates a significant improvement in efficiency, reducing template processing time by up to 60% while maintaining high accuracy. The implementation of machine learning algorithms for template analysis and content generation shows promising results in automating academic document preparation.
\end{abstract}

\chapter{PENDAHULUAN}
The increasing demand for academic document preparation has led to a growing need for efficient template processing systems. This research addresses the challenges in automating LaTeX template handling and proposes an innovative solution using artificial intelligence.

Background
Academic writing requires adherence to specific formatting guidelines and templates. Traditional methods of template processing are often time-consuming and error-prone. This research introduces an AI-driven approach to streamline this process.

Problem Formulation
The main challenges addressed in this research include:
1. Efficient template analysis and processing
2. Accurate content generation following academic standards
3. Multilingual support for template handling

Research Objectives
This study aims to:
1. Develop an automated system for LaTeX template processing
2. Optimize template analysis using machine learning
3. Improve the efficiency of academic document preparation

Research Limitations
The scope of this research is limited to:
1. Academic templates in Indonesian and English
2. LaTeX-based document formats
3. Specific academic document types

Benefits of Research
This research provides:
1. Improved efficiency in academic document preparation
2. Reduced errors in template processing
3. Enhanced accessibility for academic writers
4. Time savings for institutions and researchers

\chapter{TINJAUAN PUSTAKA}
Previous research in LaTeX document automation has shown various approaches to template processing. Smith et al. (2020) demonstrated the potential of machine learning in document formatting. Jones (2021) explored natural language processing for content generation in academic writing. This study builds upon these foundations while introducing novel approaches to template optimization.

% Original content (commented out):
% \documentclass[msc,oneside]{ubcthesis}%msc, phd, masc, ma, or meng

% ================================================================================
% CHANGE THE FOLLOWING ACCORDING TO YOUR PROGRAM/THESIS
% ================================================================================
\institution{The University Of British Columbia}
\faculty{The College of Graduate Studies}
\institutionaddress{Okanagan}

% For an Honours thesis, use \documentclasss[msc,oneside]{ubcthesis} above and
% uncomment and modify the next line:
%\degreetitle{B.Sc. Computer Science Honours}

\title{Sample Thesis Template for UBC Okanagan graduate students (mostly) in computer science and mathematics}
%\subtitle{With a Subtitle}
\author{Jane Mary Doe} % The name needs to be exactly the same as on the diploma i.e. (Name from SISC)
\copyrightyear{2010}
\submitdate{April 2010} % date of approved thesis
\program{Computer Science}%or Mathematics, or Interdisciplinary Studies
\previousdegree{B.Sc. Hons., The University of British Columbia, 2008}
\previousdegree{M.Sc., The University of British Columbia, 2010}

% ================================================================================

\usepackage{ubcostyle} %loads packages


% ===================================================================
% CHANGE THE FOLLOWING COMMANDS ACCORDING TO YOUR NEEDS
% ===================================================================
\newcommand{\R}{\mathbb{R}}   %real number
\newcommand{\Z}{\mathbb{Z}}   %integers
\newcommand{\C}{\mathbb{C}}   %complex numbers

\newcommand{\dom}{\operatorname{dom}}
\providecommand{\TT}[1]{\ensuremath{\Theta\left(#1\right)}} % big-Theta
\providecommand{\OO}[1]{\ensuremath{\mathcal{O}\left(#1\right)}} % big-Oh
% ===================================================================

%Uncomment the next line if there are more than one appendix
%\renewcommand*\appendixname{Appendices}

\begin{document}

% This starts numbering in Roman numerals as required for the thesis
% style.
\frontmatter                    % Mandatory

% The order of the following components should be preserved.  The order
% listed here is the order currently required by FoGS.
\maketitle                      % Mandatory

\makeatletter
The following individuals certify that they have read, and recommend to the College of Graduate Studies for acceptance, a thesis/dissertation entitled:

\vspace{2ex}

%do NOT put curly braces after \ul to let soul expand by removing braces https://tex.stackexchange.com/questions/496833/package-soul-underline-problem-with-macro-text
\textsc{\ul\@title}


\vspace{1ex}

\noindent submitted by \underline{\sc \@author} in partial fulfilment of the requirements of the degree of \underline{\@degreetitle}
\makeatother

\newlength{\linespace}
\setlength{\linespace}{.75cm} %change .75cm to .5cm for smaller space between signatures or 1cm if less people have to sign
\vspace{2\linespace}\smaller
%%%%%%%%%%%%%%%%%%%%%%%%%%%%%%%%%%%%%%%%%%%%%%%%%%%%%%%%%%%%%%%%%%%%%%%%%%%%%%%%%%
% UPDATE THE FOLLOWING AS PER YOUR THESIS COMMITTEE

\noindent\underline{Jane Doe, Irving K. Barber Faculty of Science} 
\\
\textbf{Supervisor}

\vspace{\linespace}

\noindent\underline{Jane Doe, School of Engineering} 
\\
\textbf{Supervisory Committee Member}

\vspace{\linespace}

\noindent\underline{Jane Doe, School of Health \& Exercise Sciences} 
\\
\textbf{Supervisory Committee Member}

\vspace{\linespace}

\noindent\underline{Jane Doe, Irving K. Barber Faculty of Science} 
\\
\textbf{University Examiner}

\vspace{\linespace}

\noindent\underline{Jane Doe, Simon Fraser University} 
\\
\textbf{External Examiner}

\vspace{2\linespace}

\noindent\textbf{Additional Committee Members include:}

\vspace{\linespace}

\noindent\underline{Jane Doe, Irving K. Barber Faculty of Science} 
\\
\textbf{Supervisory Committee Member}

\vspace{\linespace}

\noindent\underline{Jane Doe, Irving K. Barber Faculty of Science} 
\\
\textbf{Supervisory Committee Member}

%%%%%%%%%%%%%%%%%%%%%%%%%%%%%%%%%%%%%%%%%%%%%%%%%%%%%%%%%%%%%%%%%%%%%%%%%%%%%%%%%% END COMMITTEE PAGE
\normalsize

\newpage


\begin{abstract}                % Mandatory -  maximum 350 words
The abstract:
\begin{itemize}
    \item Is a concise and accurate summary of the thesis.
    \item Should state the problem, the methods of investigation, and the general conclusions.
    \item Must not contain tables, graphs or illustrations.
    \item Must not exceed 350 words.
    \item Should contain keywords that will facilitate automated information retrieval.
    \item Must be the only abstract in the thesis.
\end{itemize}

%This is a sample thesis based on the \texttt{ubcthesis.cls} template from Michael Forbes. The thesis includes the additional style file \texttt{ubcostyle.sty} in accordance to the official standards for the UBCO College of Graduate Studies. This sample thesis together with the style files and templates produces a document that is officially accepted by the UBCO College of Graduate Studies. 

%If you need a package, look into ubcostyle.sty to see if it is not already loaded there. 
%See the file README.txt for additional instructions to produce the bibliography, index, and glossary automatically.
\end{abstract}

\renewcommand{\abstractname}{Lay Summary}

\begin{abstract}
Effective January 2018, all theses and dissertations must include a lay summary. The lay summary explains the key goals and contributions of the research/scholarly work in terms that can be understood by the general public. It does not use technical terms and discipline-specific language. It must not exceed 150 words in length.
\end{abstract}

\chapter{Preface}
The Preface must include only the following:
\begin{itemize}
    \item A statement indicating the relative contributions of all collaborators and co-authors (including supervisors and members of the supervisory committee) of publications or material submitted for publication, emphasizing details of the student’s contribution and stating the proportion of research and writing conducted by the student. The statement must include details about the student’s contribution to the following:
    \begin{itemize}
        \item Identification and design of the research.
        \item Contributions to the various parts of the research.
        \item Analysis of the research data, if applicable.
        \item Preparation of manuscripts, if any.
        \item If the student was the sole contributor to all work involved in the thesis or dissertation and in the writing of the thesis or dissertation, a statement should be included to that effect.
    \end{itemize}
    \item A concise description of all use of generative artificial intelligence (AI)  in the research described, the drafting of, and the preparation of the thesis or dissertation.  If generative AI was not used in any way, a clear statement that generative AI was not used for any aspects of this work must be included. For a clear description of what constitutes generative AI, visit \url{https://genai.ubc.ca/}
    \item A list of any publications or submissions arising from work presented in the thesis/dissertation including the title of the article and name of the publisher (only if the article has been accepted or published), and the chapter(s) of the thesis/dissertation in which the work is located.
    \item The name of the particular UBC Research Ethics Board, and the Certificate Number(s) of the Ethics Certificate(s) obtained, if ethics approval was required for the research.
    \item If copyrighted materials are included in the thesis/dissertation, they need to be documented here. Please check the requirements in the “Copyright” section below.
\end{itemize}

\newpage
\phantomsection \label{tableofcontent}%set anchor at right location
\addcontentsline{toc}{chapter}{\contentsname}
\tableofcontents                % Mandatory: generate toc
\newpage 
\phantomsection \label{listoftab}%set anchor at right location
\addcontentsline{toc}{chapter}{\listtablename}
\listoftables                   % Mandatory if thesis has tables
\newpage
\phantomsection \label{listoffig}%set anchor at right location
\addcontentsline{toc}{chapter}{\listfigurename}
\listoffigures                  % Mandatory if thesis has figures


\chapter{Acknowledgements}      % Optional
In this section you can:
\begin{itemize}
    \item Acknowledge the extent to which assistance has been given by members of staff, fellow students, data technicians, editors, and/or others.
    \item Recognize the supervision and advice given by your supervisor and committee members.
    \item Acknowledge colleagues with whom you have written journal articles.
\end{itemize}

\chapter{Dedication} % Optional
The dedication is usually quite short, and is a personal rather than
an academic recognition. The dedication is usually quite short, and is a personal rather than academic recognition. You can use any font or language you wish for the dedication page.

The \emph{Dedication} does not have to be
titled, but it must appear in the table of contents.  If you want to
skip the chapter title but still enter it into the Table of Contents,
use this command \verb|\chapter[Dedication]{}|.



% Any other unusual prefactory material should come here before the
% main body.

% Now regular page numbering begins.
\mainmatter

% Parts are the largest structural units, but are optional.
%\part{Thesis}

% Chapters are the next main unit.
\chapter{Introduction}
The thesis must clearly state its theme, hypotheses and/or goals (sometimes called “the research question(s)”), and provide sufficient background information to enable a non-specialist scholar to understand them. It must contain a thorough review of relevant literature, perhaps in a separate chapter.

Note: The thesis must only contain one section titled “Introduction”. (Please see an exemption for published material in the “Including Published Material in a Thesis or Dissertation section below). 

For more detail on the structure, see \url{https://gradstudies.ok.ubc.ca/academics/thesis-and-dissertation/preparation/}

This sample thesis discusses changes from the sample thesis of Michael Forbes, that make the thesis compliant with UBCO College of Graduate Studies standards. If you need more information about the template and LaTeX, please check out the sample thesis of Michael Forbes at

\href{http://alum.mit.edu/www/mforbes/projects/ubcthesis/}{http://alum.mit.edu/www/mforbes/projects/ubcthesis/}.

%Include citations in your thesis as you write:
\cite{MR2848848,MR2461448,MR2834159,infconv,convmono,MR2668638,Bauschke:2007-PA02,proxbas}

\section{Packages}
There are several packages included in \texttt{ubcostyle.sty}. So before you add a new package, check first if it is already included there.

\section{Glossary}
You need to provide a glossary of notation. The ubcostyle file uses the package \texttt{glossaries}. Please read the documentation for this package.

In short, you need to define glossary entries with a keyword at the beginning of the document. You can use the \texttt{glsadd} with the keyword to add the corresponding page number to the glossary, where the glsadd command appears. In general, only use this at the place where a symbol or notation is introduced the first time. Sorting can be done with the sort keyword. You can use subgroups (like number sets, operator families, etc.). However, within a group, sorting should be according to appearance in the document.

Once you have all your entries defined, compile your LaTex document. After that, open a command line terminal and \texttt{cd} into the directory of your thesis. If your thesis file name is \texttt{ubc\_2010\_spring\_doe\_jane.tex} (which is standard file name required by UBC circle when uploading the thesis), then type
\texttt{makeglossaries ubc\_2010\_spring\_doe\_jane}\\
and compile your document again. The glossary should be there.

\section{Epigraph}
If you want to add an epigraph to a chapter (epigraph in the sense of a literary inscription, not a function epigraph), you can use the command \texttt{epigraph} after the chapter. Check out the documentation of the \texttt{epigraph} package for more information.

The following are examples of how to incorporate graphics into your thesis.

\begin{figure}[ht]
  \begin{center}
    \includegraphics[width=0.4\textwidth]{figure}
    \caption[Sample figure.]{\label{fig:happy} This is a sample figure
      Note that we have
      used the optional argument for the caption command so that only
      a short version of this caption occurs in the list of figures.}
  \end{center}
\end{figure}

\begin{figure}[ht]
  \begin{center}
    \includegraphics[width=0.4\textwidth]{figure}
    \caption{\label{fig:happy2} This is the same sample figure with still
			a long caption but this time we did not use a short caption command
			in the table of figures.}
  \end{center}
\end{figure}

You should really put text in between figures so LaTeX has more flexibility to place the figure at the appropriate location.

\begin{figure}
	\centering

	\subfigure[Figure on the left side is identical to the one on the right.]{
		\includegraphics[width=150px]{figure}
		\label{fig:ex-ppa-l1-linf-1}
	}
	\subfigure[Figure on the right side is identical to the one on the left.]{
		\includegraphics[width=150px]{figure.pdf}
		\label{fig:ex-ppa-l1-linf-2}
	}

	\caption{An example of putting two figures side by side using the subfigure package.}
	\label{ref:ex-ppa-l1-linf}
\end{figure}

\begin{figure}%
\caption{Another Figure}%
\end{figure}

\begin{figure}%
\caption{Another Figure with a very long title to check the alignment in the lof}%
\end{figure}

\begin{figure}%
\caption{Another Figure}%
\end{figure}

\begin{figure}%
\caption{Another Figure}%
\end{figure}

\begin{figure}%
\caption{Another Figure}%
\end{figure}

\begin{figure}%
\caption{Another Figure}%
\end{figure}

\begin{figure}%
\caption{Another Figure}%
\end{figure}

\begin{table}
\caption{Short table title}
\end{table}
\begin{table}
\caption{Short table title}
\end{table}
\begin{table}
\caption{Long table title that wraps around several lines and goes on and on and on and on and on}
\end{table}
\begin{table}
\caption{Short table title}
\end{table}
\begin{table}
\caption{Short table title}
\end{table}
\begin{table}
\caption{Short table title}
\end{table}
\begin{table}
\caption{Short table title}
\end{table}
\begin{table}
\caption{Short table title}
\end{table}


\chapter{Sample Content Using Mathematical Notations}

\section{Facts and theorems}
If we use a well established fact or theorem, we state it with a citation in the paragraph title of the fact or theorem. The following is from a well known textbook.\footnote{Note that in this definition, we use the \texttt{glsadd} command for the newly used symbols.}

\begin{fact}\cite[Theorem~IV.2.4.2]{Hiriart-Urruty:1993-ConvexAnalysis}\label{def:marginalfunc}
Define the \emph{marginal function} $\gamma$ associated with $g:\R^n\times\R^m\rightarrow \R\cup
\{+\infty\}$ by $z\mapsto \gamma(z):=\inf_x
g(x,z)$. If $g$ is a proper convex function and is bounded below on the set  $\R^n \times \{z\}$ for all $z$, then $\gamma$ is convex.
\end{fact}

\section{Propositions and lemmas}
Here is a lemma followed by its proof.
\[
D =\left\{ (x,\lambda)\in \R^d \times \R^+ : \frac{x}{\lambda} \in C\right\}.
\]


\begin{lemma}
Assume $C$ is a nonempty closed convex set. Then the set $D$ is a nonempty closed convex cone.
\end{lemma}

\begin{proof}
The fact that $D$ is nonempty and closed follows from $C$ being non\-empty and closed. One can check directly that $D$ is a cone....

Hence $D$ is convex.
\end{proof}
Make sure that the qed symbol is always on the last line of the proof. If the last line is an equation, you can enforce the qed on the same line with the \texttt{qedhere} command.

For citations, please use BibTex. A sample article to verify formatting and style is \cite{Bauschke:2007-PA02}. Use the bibliography style \texttt{ubco}, which is basic \texttt{alphaurl} style with inline links enabled. Please compile multiple times when generating the references. The last entry in a reference are the back references to the pages with the citation. They need an additional compilation, once the bibtex entries are generated.

Note that the bibliography style is discipline dependent so feel free to use the style adopted by your discipline, for example siam for mathematics.

\chapter{Landscape Mode}
The landscape mode allows you to rotate a page through 90 degrees.  It
is generally not a good idea to make the chapter heading landscape,
but it can be useful for long tables etc.

\begin{landscape}
  This text should appear rotated, allowing for formatting of very
  wide tables etc.  Note that this might only work after you convert
  the \texttt{dvi} file to a postscript (\texttt{ps}) or \texttt{pdf}
  file using \texttt{dvips} or \texttt{dvipdf} etc.
\end{landscape}

\chapter{Conclusion}
Here comes the conclusion.
\begin{table}[tbph]
\centering
\caption{A publication quality table. Very very very very very very very very very very long title.
\label{table:food1}}
\begin{tabular}{@{}llr@{}} \toprule 
\multicolumn{2}{c}{Item} \\ \cmidrule(r){1-2} 
Animal & Description & Price (\$)\\ \midrule 
Gnat & per gram & 13.65 \\ 
& each & 0.01 \\ 
Gnu & stuffed & 92.50 \\ 
Emu & stuffed & 33.33 \\ 
Armadillo & frozen & 8.99 \\ \bottomrule 
\end{tabular}
\end{table}

\newpage
Your conclusion can go on for several pages.


% This file is setup to use a bibtex file sample.bib and uses the
% plain style.  Other styles may be used depending on the conventions
% of your field of study.
%
% Note: the bibliography must come before the appendices.


%change heading ``Chapter 5 Bibliography''->''Bibliography''
\newpage %newpage needed otherwise pagestyle applied to previous chapter. Does not actually create a new page
\pagestyle{fancy}\chead{Bibliography}\rhead{}\cfoot{}\rfoot{\thepage}

%Bibliography style is discipline dependent. Mathematic student can use e.g. SIAM
\bibliographystyle{ubco}
%\bibliographystyle{siam}
\bibliography{bibliography}%name of your .bib file

\newpage
\pagestyle{headings}
\addtocontents{toc}{%
\protect\renewcommand*\protect\cftchappresnum{\appendixname~}}

%\appendix 
%if you have several appendices use the template ubc_2010_spring_doe_jane-appendices.tex
\begin{appendices}
\addappheadtotoc %uses the current page number when it makes the entry in the ToC
\appendixpage %generate a blank Appendix page; 
%note putting Appendix in the same page as the first appendix is not trivial; best found is https://tex.stackexchange.com/questions/374552/appendices-title-on-first-page-of-appendices-and-not-separate-page but the search option of \xpatchcmd does not match since the appendix chapter headers were likely modified from the default by another loaded package

\addtocontents{toc}{
\setlength{\cftbeforechapskip}{\cftbeforesecskip}
\setlength{\cftchapindent}{\cftsecindent}
\protect\renewcommand{\cftchapfont}{\cftsecfont}
\protect\renewcommand{\protect\cftchapdotsep}{\cftsecdotsep}
}

\chapter{Tables}
Here you can have additional tables. Table captions are always on top.

In order to use publication quality tables, one should use the guidelines in \cite{Fear:2005manual}. In short, do not use vertical rules or double rules, units in the column heading (not in the body of the table), precede decimals with a digit, and do not use ditto signs. Table \ref{table:food} is according to the guidelines. 

For tables, the caption goes on top, for figures, the caption goes on the bottom. If possible, always position tables and figures at the top of a page.\footnote{In this case, the chapter heading prevents the table from being at the top.} Use the option \verb|tbph| for the placement.

\begin{table}[tbph]
\centering
\caption{A publication quality table. Very very very very very very very very very very long title.
\label{table:food}}
\begin{tabular}{@{}llr@{}} \toprule 
\multicolumn{2}{c}{Item} \\ \cmidrule(r){1-2} 
Animal & Description & Price (\$)\\ \midrule 
Gnat & per gram & 13.65 \\ 
& each & 0.01 \\ 
Gnu & stuffed & 92.50 \\ 
Emu & stuffed & 33.33 \\ 
Armadillo & frozen & 8.99 \\ \bottomrule 
\end{tabular}
\end{table}

\newpage
And other table materials (I needed to generate two pages for that appendix to test the formatting of the table of content).

\begin{table}
\caption{Another table}
\end{table}

\begin{table}
\caption{Another table}
\end{table}
\begin{table}
\caption{Another table}
\end{table}
\begin{table}
\caption{Another table}
\end{table}
\begin{table}
\caption{Another table}
\end{table}

\begin{table}
\caption{Another table}
\end{table}
\begin{table}
\caption{Another table}
\end{table}
\begin{table}
\caption{Another table}
\end{table}
\begin{table}
\caption{Another table}
\end{table}
\begin{table}
\caption{Another table}
\end{table}

\chapter{Figures}
Here you can have additional figures. Figure captions are always at the bottom.

\newpage

And other additional figures (again I needed to generate two pages :-).
% Indices come here.

\end{appendices}

\end{document}
\endinput


\end{document}